\section{Numerical illustration}
\label{sec:numerical}

We now illustrate the quantitative implications of the model using the calibrated parameters described in Section~\ref{sec:calibration}. The baseline parameters are $\Sharpe=0.35$, $r_f=3.70\%$, $\gamma=5$, $\eta=15$, $A_0=1$, $f_1=1.5$, and $f_2=25$.

At the baseline calibration, the model delivers an equilibrium volatility target $\sigma_d^* \approx 12.8\%$ (annualized), consistent with the volatility range for actively managed equity funds noted in Section~\ref{sec:calibration}. The incentive-compatible fee is $\phi^* \approx 3.95\%$, computed from the IC fee schedule \eqref{eq:phi-IC} and the investor's first-order condition \eqref{eq:investor-FOC}. The manager's equilibrium utility $U_M(\sigma_d^*, \phi^*) \approx 0.050 > 0$, confirming that the participation constraint $U_M \ge 0$ is slack at the calibrated optimum. The affine approximation error between the rule $\alpha + \beta\sigma_d$ and the exact IC fee $\phi(\sigma_d)$ is approximately 0.08 percentage points on average in absolute terms, evaluated within $\pm 5\%$ of $\sigma_d^*$ (see Figure~\ref{fig:affine-approx}). Figure~\ref{fig:incentive-alignment} confirms incentive alignment: the investor's utility is maximized at $\sigma_d^*$ along the IC curve, and the manager's utility under the fixed IC fee $\phi^*$ is itself maximized at exactly $\sigma_d^*$. Figure~\ref{fig:welfare} provides the welfare comparison against a flat-fee benchmark of $\phi_{\rm mkt}=2\%$ (an empirically typical management fee). At the flat fee, the manager freely chooses $\sigma^{\rm free}\approx 19.1\%>\sigma_d^*$, exploiting convex flows. Under the flat fee the investor's utility $U_I(\sigma^{\rm free},\phi_{\rm mkt})\approx -0.006$, while the IC contract achieves $U_I(\sigma_d^*,\phi^*)\approx 0.001$, a welfare gain of $\Delta U_I\approx 0.007$. Notably, the investor's participation constraint ($U_I\ge 0$) is satisfied under the IC contract but violated under the flat-fee benchmark, because the manager's over-risking at $\sigma^{\rm free}\approx 19.1\%$ generates variance large enough to make the investor's mean-variance utility negative despite the high expected return. The right panel shows this welfare gain is decreasing in $f_2$: IC enforcement delivers the largest improvement over the flat-fee benchmark at moderate convexity, and the gain diminishes as $f_2$ increases because the IC fee must fall sharply to maintain incentive alignment. At sufficiently high convexity ($f_2\gtrsim 40$), the IC contract no longer improves on the flat-fee benchmark.

Beyond the baseline, Figures~\ref{fig:comp-f2}--\ref{fig:comp-risk} illustrate how the equilibrium target $\sigma_d^*$ and the IC fee $\phi(\sigma_d^*)$ vary with the key parameters, computed at the Section~\ref{sec:calibration} calibration baseline while sweeping each parameter over empirically plausible ranges (documented in Appendix~\ref{app:numerical}). Figure~\ref{fig:flow-performance} shows the flow-performance relationship. Figure~\ref{fig:ic-fee-curve} shows the IC mapping: for each target volatility $\sigma_d$, there is a fee level that makes the manager willing to implement it; higher convexity shifts the schedule downward, reflecting the stronger incentive to over-risk when performance chasing is more pronounced. Figure~\ref{fig:opt-sigma-vs-f2} shows that stronger convexity lowers the investor's optimal target, consistent with the comparative statics: sustaining a high target becomes more contractually costly when flows are more convex. Figures~\ref{fig:comp-sharpe} and~\ref{fig:comp-risk} illustrate that higher fund skill raises the optimal target while higher investor risk aversion lowers it, and that a more risk-averse manager requires a lower IC fee to implement any given target, since greater natural risk aversion reduces the need for contractual discipline.

Throughout these exercises, the manager's participation constraint $U_M(\sigma_d^*,\phi(\sigma_d^*)) \ge 0$ and the investor's participation constraint $U_I(\sigma_d^*,\phi^*) \ge 0$ (normalized to zero reservation utility) are both satisfied throughout the calibrated parameter ranges. Notably, the investor's constraint is violated under the flat-fee benchmark --- the manager's over-risking at $\sigma^{\rm free}\approx 19.1\%$ generates a variance penalty that makes the investor's mean-variance utility negative --- whereas the IC contract restores a non-negative utility. The limited-liability constraint $\phi^* \ge 0$ holds across all parameter sweeps in Appendix~\ref{app:numerical}: the IC fee formula $\phi(\sigma_d) = 2\,\partial_\sigma \E[A] / (\eta\,\partial_\sigma \V[A])$ inherits the sign of $\partial_\sigma \E[A]$, which is strictly positive whenever $S > 0$ or $f_2 > 0$. This confirms that, under the calibration, the relevant binding constraint is incentive compatibility rather than participation.

A practical question concerns how the contract $\phi(\sigma_d)$ is implemented: the investor must specify a target volatility $\sigma_d$ and the manager must implement it. In practice, the realized volatility $\hat{\sigma}$ must be estimated from a time series of fund returns, introducing estimation error. Over a standard evaluation horizon (e.g., quarterly or annual), estimation noise is non-negligible. A natural implementation maps a window of realized returns to an estimated $\hat{\sigma}$, then sets the fee according to the affine rule $\phi \approx \alpha + \beta\hat{\sigma}$. Two complications arise. First, estimation error in $\hat{\sigma}$ may cause the contract to inadvertently reward or penalize the manager for luck in realized returns rather than the risk level chosen. Second, managers have an incentive to manage reported volatility through return smoothing or infrequent pricing of illiquid assets, as documented in \citet{GetmanskyLoMakarov2004} and \citet{BollenPool2009} for hedge funds. In equity mutual fund contexts, mark-to-market pricing and regulatory reporting requirements reduce but do not eliminate this concern. The practical solution is to evaluate the contract over sufficiently long horizons and to use realized volatility measures that are robust to smoothing, such as high-frequency range-based estimators. A full extension of the model to estimated volatility is beyond the scope of the present analysis but constitutes a natural direction for future work.

Taken together, these results highlight how the model translates observable features of investor behavior into implications for contract design. For both investors and managers, the incentive-compatible fee schedule and its affine approximation provide a transparent mapping from a volatility target to a corresponding fee level, sustaining incentive alignment under realistic market conditions.
